% Juknis Pembaruan RPS untuk Program Studi D4 Teknik Informatika
% Compile with xelatex from docs/juknis-rps/ directory.
% Format mengikuti kop surat resmi Jurusan Teknologi Informasi, sebagaimana
% dicontohkan pada docs/JUKNIS YUDISIUM.pdf.
\documentclass[12pt,a4paper]{report}

\newcommand{\LabName}{Pembaruan RPS}

% Shared preamble for all Pedoman Laboratorium documents.
% Each lab file must define \LabName before \input-ing this file.
% Compile with xelatex from the repository root so assets/ resolves correctly.

\usepackage{fontspec}
\usepackage{polyglossia}
\setdefaultlanguage{indonesian}
\usepackage[margin=2.4cm, left=2.8cm, top=2.3cm, bottom=2.2cm,
            headheight=64pt, headsep=8pt, footskip=20pt]{geometry}
\usepackage{graphicx}
\usepackage[table]{xcolor}
\usepackage{tabularray}
\usepackage{enumitem}
\usepackage{fancyhdr}
\usepackage{titlesec}
\usepackage{tikz}
\usepackage{eso-pic}
\usepackage{hyperref}
\usepackage{setspace}
\usepackage[protrusion=true,expansion=false]{microtype}
\usepackage{xstring}

\usetikzlibrary{calc,positioning,arrows.meta}
\UseTblrLibrary{varwidth}

% ----- Fonts ---------------------------------------------------------------
\defaultfontfeatures{Ligatures=TeX}
\setmainfont{Times New Roman}
\setsansfont{Arial}
\newfontfamily\titlefont{Arial}
\newfontfamily\headingfont{Cambria}

% ----- Colors --------------------------------------------------------------
\definecolor{polinemablue}{HTML}{1679BF}
\definecolor{polinemagrey}{HTML}{5B6A72}
\definecolor{headblue}{HTML}{D9EEF7}
\definecolor{headgray}{HTML}{D9D9D9}
\definecolor{headgreen}{HTML}{C6E0B4}

% ----- Asset paths (CWD = docs/pedoman-laboratorium/ during compilation) --
% logo-left/right.png live in book/src/images/; logo-polinema-transparent.png in assets/
\graphicspath{{../../book/src/images/}{../../assets/}{assets/}{./}}

% ----- Tblr templates (suppress caption noise) ----------------------------
\DeclareTblrTemplate{conthead-text}{empty}{}
\DeclareTblrTemplate{contfoot-text}{empty}{}
\SetTblrTemplate{caption-tag}{empty}
\SetTblrTemplate{caption-sep}{empty}
\SetTblrTemplate{caption}{empty}
\SetTblrTemplate{capcont}{empty}
\SetTblrTemplate{conthead}{empty}
\SetTblrTemplate{conthead-text}{empty}
\SetTblrTemplate{contfoot}{empty}
\SetTblrTemplate{contfoot-text}{empty}
\SetTblrTemplate{caption-lot}{empty}

% ----- Name helper: avoid double periods when a name already ends in one --
% (e.g. degree abbreviations like "M.T." or "Ph.D.") when used at the end
% of a sentence that appends its own period.
\newcommand{\SentencePeriod}[1]{\IfEndWith{#1}{.}{}{.}}

% ----- Ketua Jurusan (department head) constants ---------------------------
% Shared across all labs; second signatory on \PedomanPengesahan.
\newcommand{\KetuaJurusan}{Mungki Astiningrum, S.T., M.Kom.}
\newcommand{\KetuaJurusanNIP}{19771030 200501 2 001}

% ----- Page chrome (blue rules on every page) ------------------------------
\newcommand{\PageChrome}{%
  \begin{tikzpicture}[remember picture,overlay]
    \draw[line width=1.2pt, color=polinemablue]
      ($(current page.north west)+(1.0cm,-1.25cm)$) --
      ($(current page.north east)+(-1.0cm,-1.25cm)$);
    \draw[line width=1.2pt, color=polinemablue]
      ($(current page.south west)+(1.0cm,1.1cm)$) --
      ($(current page.south east)+(-1.0cm,1.1cm)$);
  \end{tikzpicture}%
}
\AddToShipoutPictureBG{\PageChrome}

% ----- Header/footer -------------------------------------------------------
\pagestyle{fancy}
\fancyhf{}
\renewcommand{\headrulewidth}{0pt}
\renewcommand{\footrulewidth}{0pt}
\fancyhead[L]{\includegraphics[height=42pt,keepaspectratio]{logo-left.png}}
\fancyhead[R]{\includegraphics[height=42pt,keepaspectratio]{logo-right.png}}
\fancyfoot[R]{\textcolor{polinemagrey}{\headingfont\itshape\fontsize{9pt}{9pt}\selectfont
  D4 Teknik Informatika\ \textbar\ POLINEMA\ \ \thepage}}

\fancypagestyle{plain}{%
  \fancyhf{}
  \fancyhead[L]{\includegraphics[height=42pt,keepaspectratio]{logo-left.png}}
  \fancyhead[R]{\includegraphics[height=42pt,keepaspectratio]{logo-right.png}}
  \fancyfoot[R]{\textcolor{polinemagrey}{\headingfont\itshape\fontsize{9pt}{9pt}\selectfont
    D4 Teknik Informatika\ \textbar\ POLINEMA\ \ \thepage}}
  \renewcommand{\headrulewidth}{0pt}
  \renewcommand{\footrulewidth}{0pt}
}

% ----- Chapter / section headings ------------------------------------------
\titleformat{\chapter}[display]
  {\centering\headingfont\bfseries\Large}
  {BAB \thechapter}
  {0.35em}
  {\MakeUppercase}
\titlespacing*{\chapter}{0pt}{-1.2em}{0.7em}

\titleformat{\section}
  {\bfseries\normalsize}
  {\thesection}
  {0.6em}
  {}
\titlespacing*{\section}{0pt}{5pt}{1.5pt}

\titleformat{\subsection}
  {\bfseries\normalsize}
  {\thesubsection}
  {0.6em}
  {}
\titlespacing*{\subsection}{0pt}{6pt}{2pt}

% ----- Misc layout ---------------------------------------------------------
\setstretch{1.03}
\setlength{\parindent}{0pt}
\setlength{\parskip}{2pt}
\setlength{\emergencystretch}{2.5em}
\setlength{\tolerance}{800}
\hyphenation{pe-rang-kat a-ka-de-mik skrip-si re-ka-ya-sa peng-uji pem-bim-bing pe-ne-li-ti-an pe-ngem-bang-an}
\XeTeXlinebreakpenalty=0
\raggedbottom
\setlist[itemize]{leftmargin=1.4em, topsep=1pt, itemsep=0pt, parsep=0pt}
\setlist[enumerate]{leftmargin=1.8em, topsep=1pt, itemsep=0pt, parsep=0pt}

\hypersetup{colorlinks=true, linkcolor=black, urlcolor=polinemablue}

% ----- Cover page macro ----------------------------------------------------
% Uses \LabName (e.g. "SOFTWARE ENGINEERING"),
%      \LabKoor (e.g. "Habibie")
\newcommand{\PedomanCover}{%
  \begin{titlepage}
  \thispagestyle{empty}
  % Blue bottom band
  \begin{tikzpicture}[remember picture,overlay]
    \fill[polinemablue]
      (current page.south west) rectangle ([yshift=4.2cm]current page.south east);
    \node[white, font=\titlefont\bfseries]
      at ([yshift=3.4cm]current page.south)
      {\fontsize{14pt}{17pt}\selectfont JURUSAN TEKNOLOGI INFORMASI};
    \node[white, font=\titlefont\bfseries]
      at ([yshift=2.7cm]current page.south)
      {\fontsize{14pt}{17pt}\selectfont POLITEKNIK NEGERI MALANG};
    \node[white, font=\titlefont\bfseries]
      at ([yshift=1.9cm]current page.south)
      {\fontsize{15pt}{18pt}\selectfont 2025};
  \end{tikzpicture}

  \vspace*{1.8cm}

  \begin{center}
    {\titlefont\bfseries\fontsize{17pt}{20pt}\selectfont
      BUKU PEDOMAN LABORATORIUM\par}
    \vspace{0.35cm}
    {\color{polinemablue}\rule{0.55\textwidth}{1.5pt}\par}
    \vspace{0.45cm}
    {\titlefont\bfseries\fontsize{16pt}{18pt}\selectfont
      \textcolor{polinemablue}{\MakeUppercase{\LabName}}\par}
  \vspace{0.25cm}

  \end{center}

  \vspace{1.0cm}

  \begin{center}
    \includegraphics[height=4cm,keepaspectratio]{logo-right.png}
  \end{center}

  \vspace{0.6cm}

  \begin{center}
    {\headingfont\bfseries\fontsize{24pt}{28pt}\selectfont
      \textcolor{polinemablue}{POLINEMA}\par}
  \end{center}

  \end{titlepage}
}

% ----- Daftar Isi title ----------------------------------------------------
\renewcommand{\contentsname}{DAFTAR ISI}

% ----- Organizational responsibility-flow diagram ---------------------------
% Generic, lab-agnostic (no \LabKoor / person names) -- used in Chapter 2 of
% every lab file in place of the old inline "$\rightarrow$" sentence.
\tikzset{
  orgtop/.style={draw, rounded corners, fill=polinemablue, text=white,
    font=\bfseries\fontsize{8}{9}\selectfont, align=center,
    text width=6.0cm, minimum height=0.85cm, inner sep=4pt},
  orgmid/.style={draw, rounded corners, fill=headblue,
    font=\bfseries\fontsize{8}{9}\selectfont, align=center,
    text width=6.0cm, minimum height=0.85cm, inner sep=4pt},
  orgleaf/.style={draw, rounded corners, fill=headgray,
    font=\fontsize{7.5}{8.5}\selectfont, align=center,
    text width=2.9cm, minimum height=1.05cm, inner sep=3pt},
}
\newcommand{\OrgChart}{%
\begin{center}
\begin{tikzpicture}[
    >={Stealth[length=2mm]}, thick,
    node distance=9mm and 4mm,
  ]
  \node[orgtop] (jurusan) {Ketua Jurusan\\Teknologi Informasi};
  \node[orgmid, below=of jurusan] (prodi) {Ketua Program Studi\\D4 Teknik Informatika};
  \node[orgmid, below=of prodi] (kalab) {Kepala Laboratorium};

  \node[orgleaf, below left=13mm and 4.5cm of kalab] (teknisi) {Teknisi /\\Laboran};
  \node[orgleaf, right=4mm of teknisi] (pembimbing) {Dosen\\Pembimbing};
  \node[orgleaf, right=4mm of pembimbing] (penguji) {Dosen\\Penguji};
  \node[orgleaf, right=4mm of penguji] (mentor) {Mentor Industri /\\Asisten Capstone};

  \draw[->] (jurusan) -- (prodi);
  \draw[->] (prodi) -- (kalab);
  \draw[->] (kalab.south) -- ++(0,-6mm) -| (teknisi.north);
  \draw[->] (kalab.south) -- ++(0,-6mm) -| (pembimbing.north);
  \draw[->] (kalab.south) -- ++(0,-6mm) -| (penguji.north);
  \draw[->] (kalab.south) -- ++(0,-6mm) -| (mentor.north);
\end{tikzpicture}
\end{center}%
}

% ----- Lembar Pengesahan (approval page) macro ------------------------------
% Uses \LabName, \LabKoor, \LabKoorNIP (per-lab macros defined before
% \input{preamble}) plus the shared \KetuaJurusan / \KetuaJurusanNIP above.
\newcommand{\PedomanPengesahan}{%
% ===========================================================================
\chapter*{Lembar Pengesahan}
\addcontentsline{toc}{chapter}{Lembar Pengesahan}
% ===========================================================================

Buku Pedoman Laboratorium \emph{\LabName} ini telah diperiksa dan disahkan sebagai acuan resmi penyelenggaraan capstone project dan skripsi di Program Studi D4 Teknik Informatika, Jurusan Teknologi Informasi, Politeknik Negeri Malang.

\vspace{1cm}

\begin{tblr}{width=\textwidth,colspec={X[c,m]X[c,m]},rowsep=3pt,colsep=3pt}
Malang, \underline{\hspace{3.5cm}} & \\[4pt]
Kepala Laboratorium & Mengesahkan, \\
\emph{\LabName} & Ketua Jurusan Teknologi Informasi \\[28mm]
\underline{\hspace{5.5cm}} & \underline{\hspace{5.5cm}} \\
\textbf{\LabKoor} & \textbf{\KetuaJurusan} \\
\LabKoorNIP & NIP. \KetuaJurusanNIP \\
\end{tblr}
}

\usepackage{amssymb}
\usepackage{array}
\newcolumntype{L}[1]{>{\raggedright\arraybackslash}p{#1}}
\newcolumntype{C}[1]{>{\centering\arraybackslash}p{#1}}

% ----- No decorative page-rule chrome; the kop surat below is the header ---
\ClearShipoutPictureBG{}

% ----- Section headings: lettered A, B, C ... like the reference juknis ----
\renewcommand{\thesection}{\Alph{section}}
\titleformat{\section}
  {\bfseries\normalsize}
  {\thesection.}
  {0.6em}
  {}
\titlespacing*{\section}{0pt}{12pt}{4pt}

% ----- Nested list labels a., b., c. (matches reference's 1. -> a. style) --
\setlist[enumerate,2]{label=\alph*., leftmargin=1.6em}

% ----- Body paragraph indent (reference indents first lines of prose) ------
\setlength{\parindent}{0.75cm}

% ----- Kop surat (official letterhead), repeated on every page -------------
\newcommand{\KopSurat}{%
  \noindent
  \begin{minipage}[t]{0.15\textwidth}
    \vspace{2pt}\centering\includegraphics[height=1.7cm,keepaspectratio]{logo-right.png}
  \end{minipage}%
  \begin{minipage}[t]{0.85\textwidth}
    \setlength{\parskip}{0pt}\centering
    {\bfseries\fontsize{13pt}{15pt}\selectfont KEMENTERIAN PENDIDIKAN TINGGI, SAINS,\\
      DAN TEKNOLOGI\\[1pt]}
    {\bfseries\fontsize{13pt}{15pt}\selectfont POLITEKNIK NEGERI MALANG\\[1pt]}
    {\bfseries\fontsize{12pt}{14pt}\selectfont JURUSAN TEKNOLOGI INFORMASI\\[1pt]}
    {\bfseries\fontsize{9pt}{11pt}\selectfont Jl. Soekarno Hatta No.9 Malang 65141\\[1pt]}
    {\bfseries\fontsize{9pt}{11pt}\selectfont Telp (0341) 404424 -- 404425 Fax (0341) 404420\\[1pt]}
    {\bfseries\fontsize{9pt}{11pt}\selectfont
      Laman:/\underline{\textcolor{polinemablue}{www.polinema.ac.id}}\ Email:cs@polinema.ac.id}
  \end{minipage}%
}

\geometry{headheight=110pt, headsep=8pt, top=4.6cm}
\fancyhf{}
\fancyhead[C]{\KopSurat}
\fancyfoot[R]{\thepage}
\renewcommand{\headrulewidth}{1.2pt}
\renewcommand{\footrulewidth}{0pt}
\fancypagestyle{plain}{%
  \fancyhf{}
  \fancyhead[C]{\KopSurat}
  \fancyfoot[R]{\thepage}
  \renewcommand{\headrulewidth}{1.2pt}
  \renewcommand{\footrulewidth}{0pt}
}

% ----- Flow-diagram boxes (local to this document) --------------------------
\usetikzlibrary{shapes.geometric}
\tikzset{
  flowbox/.style={draw, rounded corners, fill=headblue,
    font=\bfseries\fontsize{8.5}{10}\selectfont, align=center,
    text width=6.4cm, minimum height=1.05cm, inner sep=5pt},
  flowdecision/.style={draw, diamond, aspect=2.6, fill=headgray,
    font=\bfseries\fontsize{8}{9.5}\selectfont, align=center,
    text width=3.2cm, inner sep=2pt},
  flowend/.style={draw, rounded corners, fill=headgreen,
    font=\bfseries\fontsize{8.5}{10}\selectfont, align=center,
    text width=6.4cm, minimum height=1.05cm, inner sep=5pt},
}
\newcommand{\AlurRPS}{%
\begin{center}
\begin{tikzpicture}[>={Stealth[length=2mm]}, thick, node distance=8mm]
  \node[flowbox] (langkah1) {1. Dosen Pengembang RPS menyusun/merevisi RPS};
  \node[flowbox, below=of langkah1] (langkah2) {2. Koordinator RMK/KBK memeriksa RPS menggunakan Checklist Validasi (Bagian G angka 3)};
  \node[flowdecision, below=13mm of langkah2] (cek) {Checklist terpenuhi?};
  \node[flowbox, below=13mm of cek] (langkah4) {4. Ka Prodi mengesahkan RPS (tanda tangan blok OTORISASI)};
  \node[flowend, below=of langkah4] (langkah5) {5. RPS diunggah ke SIAKAD; pembaruan dicatat pada log revisi};

  \draw[->] (langkah1) -- (langkah2);
  \draw[->] (langkah2) -- (cek);
  \draw[->] (cek) -- node[right, font=\small]{Ya} (langkah4);
  \draw[->] (langkah4) -- (langkah5);
  \draw[->] (cek.west) -- ++(-2.6cm,0) node[midway, above, font=\small]{Tidak} |- (langkah1.west);
\end{tikzpicture}
\end{center}%
}

\begin{document}

% ===========================================================================
\begin{center}
  {\bfseries\fontsize{14pt}{17pt}\selectfont PETUNJUK TEKNIS\par}
  \vspace{2pt}
  {\bfseries\fontsize{14pt}{17pt}\selectfont PEMBARUAN RENCANA PEMBELAJARAN SEMESTER (RPS)\par}
  \vspace{2pt}
  {\fontsize{11pt}{13pt}\selectfont Program Studi D4 Teknik Informatika\par}
\end{center}

\vspace{0.6cm}

% ===========================================================================
\section{Tujuan}
% ===========================================================================

Petunjuk teknis ini menjamin setiap RPS Program Studi D4 Teknik Informatika tetap mutakhir dan melalui siklus review, validasi, dan pengesahan yang formal dan konsisten. Ketentuan ini disusun untuk menjawab temuan evaluasi akreditasi internasional IABEE pada Kriteria 2.1, yang mencatat bahwa Program Studi belum memiliki mekanisme formal yang memastikan seluruh RPS direview, divalidasi, dan disahkan secara konsisten. Yang dimaksud evaluator bukan sekadar ketiadaan hasil akhir seperti RPS yang sudah ditandatangani, melainkan ketiadaan prosedur yang terdokumentasi: siapa yang memeriksa, kapan pemeriksaan dilakukan, kriteria apa yang digunakan, dan bagaimana jejaknya dicatat.

% ===========================================================================
\section{Ruang Lingkup}
% ===========================================================================

Petunjuk teknis ini berlaku untuk seluruh RPS mata kuliah Kurikulum 2025, yang berjumlah 59 mata kuliah, beserta seluruh Koordinator RMK/KBK di Program Studi D4 Teknik Informatika yang bertanggung jawab atas pemeriksaannya, mulai dari munculnya pemicu pembaruan RPS sampai dengan pengesahan, pengunggahan ke SIAKAD, dan pengarsipan log revisi.

% ===========================================================================
\section{Dokumen Acuan}
% ===========================================================================

\begin{enumerate}
\item Peraturan Menteri Pendidikan Tinggi, Sains, dan Teknologi Nomor 39 Tahun 2025 tentang Penjaminan Mutu Pendidikan Tinggi.
\item Nota Dinas Wakil Direktur I Nomor 19/WADIR.I/KM/2025 tanggal 6 Februari 2025, tentang penerapan pendekatan \emph{Outcome-Based Education} (OBE) di seluruh program studi.
\item Nota Dinas Wakil Direktur I Nomor 11/WADIR I/KM/2026 tanggal 15 Januari 2026, tentang Penataan Kurikulum T.A.\ 2026/2027 sesuai Standar Penjaminan Mutu Pendidikan Tinggi, IKU 2026, dan Akreditasi Internasional. Nota dinas ini berlaku efektif mulai Semester Ganjil T.A.\ 2026/2027 untuk seluruh angkatan.
\item Templat RPS Program Studi D4 Teknik Informatika Kurikulum 2025, khususnya blok OTORISASI (Dosen Pengembang RPS, Koordinator RMK, Ka Prodi) dan kolom TGL.\ PENYUSUNAN.
\end{enumerate}

% ===========================================================================
\section{Prinsip Pelaksanaan}
% ===========================================================================

\begin{enumerate}
\item \textbf{Terdokumentasi dan dapat ditelusuri}, yaitu setiap review, validasi, dan pengesahan RPS dicatat dan dapat dibuktikan jejaknya.
\item \textbf{Konsisten}, yaitu pemeriksaan RPS mengikuti Checklist Validasi baku yang sama untuk seluruh mata kuliah.
\item \textbf{Sederhana}, yaitu petunjuk teknis ini memanfaatkan rantai otorisasi (Dosen Pengembang RPS, Koordinator RMK, Ka Prodi) yang sudah berjalan pada templat RPS, tanpa menambah peran atau kolom baru.
\item \textbf{Dapat dipertanggungjawabkan}, yaitu pembaruan RPS hanya dilakukan karena pemicu yang sah sebagaimana diatur pada Bagian G angka 1.
\end{enumerate}

% ===========================================================================
\section{Pihak yang Terlibat dan Bertanggung Jawab}
% ===========================================================================

\begin{center}
\begin{tblr}{width=\textwidth,colspec={|Q[l,m,4.2cm]|X[l,m]|},hlines,vlines,hspan=minimal,rowsep=4pt,colsep=4pt,row{1}={bg=headblue,font=\bfseries}}
Pihak & Tanggung Jawab \\
Dosen Pengembang RPS & Menyusun atau merevisi RPS sesuai pemicu yang berlaku. \\
Koordinator RMK/KBK & Memeriksa RPS menggunakan Checklist Validasi sebelum diajukan untuk pengesahan Ka Prodi; mengembalikan RPS kepada Dosen Pengembang apabila checklist belum terpenuhi. \\
Ka Prodi & Mengesahkan RPS dengan menandatangani blok OTORISASI setelah checklist terpenuhi. \\
\end{tblr}
\end{center}

Istilah ``RMK'' (dipakai pada berkas RPS) dan ``KBK'' (dipakai pada dokumen tanggapan akreditasi) merujuk pada peran yang sama.

% ===========================================================================
\section{Dokumen dan Data yang Disiapkan}
% ===========================================================================

\begin{enumerate}
\item Berkas RPS terkini untuk mata kuliah yang akan diperbarui.
\item Checklist Validasi RPS (lihat Bagian G angka 3).
\item Templat RPS dengan blok OTORISASI dan kolom TGL.\ PENYUSUNAN.
\item Log revisi RPS Program Studi.
\item Akses unggah RPS pada SIAKAD.
\end{enumerate}

% ===========================================================================
\section{Prosedur Pelaksanaan}
% ===========================================================================

\begin{enumerate}
\item \textbf{Pemicu Pembaruan RPS.} RPS tidak diperbarui secara sembarangan. Pembaruan hanya dilakukan apabila terjadi salah satu dari tiga kondisi berikut, agar perubahan tetap dapat dipertanggungjawabkan alasannya.
  \begin{enumerate}
  \item Siklus rutin setiap semester, yaitu peninjauan yang dilakukan menjelang Semester Ganjil maupun menjelang Semester Genap dimulai, mengikuti kalender akademik Politeknik Negeri Malang (lihat angka 6).
  \item Perubahan kebijakan resmi, yaitu terbitnya peraturan atau nota dinas baru yang berdampak pada struktur kurikulum atau capaian pembelajaran, misalnya Nota Dinas 11/WADIR I/KM/2026.
  \item Temuan evaluasi, yaitu hasil tracer study, evaluasi pembelajaran mahasiswa, atau temuan asesmen akreditasi yang merekomendasikan perubahan pada suatu mata kuliah.
  \end{enumerate}

\item \textbf{Penyusunan atau Revisi RPS.}
  \begin{enumerate}
  \item Dosen Pengembang RPS menyusun atau merevisi RPS sesuai pemicu yang berlaku.
  \end{enumerate}

\item \textbf{Pemeriksaan oleh Koordinator RMK/KBK.}
  \begin{enumerate}
  \item Koordinator RMK/KBK memeriksa lima hal berikut sebelum RPS diajukan untuk pengesahan Ka Prodi: kesesuaian CPL, CPMK, dan Sub-CPMK dalam RPS; kesesuaian bobot SKS dengan distribusi resmi Kurikulum 2025; referensi/pustaka masih relevan dan tersedia; blok OTORISASI lengkap (Dosen Pengembang RPS, Koordinator RMK, Ka Prodi); serta rubrik penilaian lengkap dan konsisten dengan bentuk penilaian pada RPS.
  \item Apabila checklist belum terpenuhi, RPS dikembalikan kepada Dosen Pengembang untuk diperbaiki, dan pemeriksaan diulang dari angka ini.
  \end{enumerate}

\item \textbf{Pengesahan oleh Ka Prodi.}
  \begin{enumerate}
  \item Apabila checklist terpenuhi, Ka Prodi mengesahkan RPS dengan menandatangani blok OTORISASI, mengikuti format yang sudah berlaku pada seluruh RPS saat ini.
  \item Kolom TGL.\ PENYUSUNAN pada RPS diperbarui menjadi tanggal revisi terakhir, tanpa menambah kolom versi baru pada templat RPS.
  \end{enumerate}

\item \textbf{Unggah SIAKAD dan Pencatatan Log Revisi.}
  \begin{enumerate}
  \item RPS yang telah disahkan diunggah ke SIAKAD.
  \item Sebagai pelengkap, Program Studi mengelola sebuah log revisi bersama yang mencatat setiap pembaruan RPS: kode dan nama mata kuliah, tanggal revisi, alasan singkat perubahan, serta nama Koordinator RMK/KBK yang memvalidasinya. Log ini menjadi bukti bahwa siklus review benar-benar berjalan dari waktu ke waktu, bukan hanya prosedur di atas kertas.
  \end{enumerate}

\item \textbf{Jadwal.}
  \begin{enumerate}
  \item Peninjauan RPS secara rutin dilakukan menjelang setiap semester dimulai, mengikuti kalender akademik Politeknik Negeri Malang yang ditetapkan Program Studi setiap tahun: menjelang Semester Ganjil pada bulan Juni sampai dengan Juli (selaras dengan berlakunya Nota Dinas 11/WADIR I/KM/2026 pada Semester Ganjil T.A.\ 2026/2027), dan menjelang Semester Genap sesuai kalender akademik yang berlaku pada tahun akademik tersebut.
  \item Peninjauan menjelang suatu semester difokuskan pada RPS mata kuliah yang akan diselenggarakan pada semester tersebut.
  \item Pemicu di luar jadwal rutin ini, yaitu perubahan kebijakan resmi atau temuan evaluasi sebagaimana disebutkan pada angka 1, dapat memulai siklus pembaruan kapan pun diperlukan tanpa menunggu periode peninjauan rutin di atas.
  \end{enumerate}
\end{enumerate}

Alur pada angka 2 sampai dengan 5 di atas digambarkan pada diagram berikut.

\AlurRPS

% ===========================================================================
\section{Keluaran}
% ===========================================================================

\begin{center}
\begin{tblr}{width=\textwidth,colspec={|Q[l,m,4.6cm]|X[l,m]|},hlines,vlines,hspan=minimal,rowsep=4pt,colsep=4pt,row{1}={bg=headblue,font=\bfseries}}
Keluaran & Fungsi \\
RPS tervalidasi dan disahkan & RPS mata kuliah yang telah diperiksa Koordinator RMK/KBK dan disahkan Ka Prodi melalui blok OTORISASI. \\
Checklist Validasi terisi & Bukti pemeriksaan kesesuaian RPS sebelum pengesahan. \\
Log revisi RPS & Catatan riwayat pembaruan RPS Program Studi dari waktu ke waktu. \\
RPS pada SIAKAD & Salinan RPS terkini yang dapat diakses oleh mahasiswa dan dosen. \\
\end{tblr}
\end{center}

% ===========================================================================
\section{Ketentuan Khusus}
% ===========================================================================

\begin{enumerate}
\item RPS tidak diperbarui secara sembarangan; pembaruan hanya dilakukan karena pemicu yang sah sebagaimana diatur pada Bagian G angka 1.
\item Kolom TGL.\ PENYUSUNAN menjadi penanda revisi terakhir; templat RPS tidak menambahkan kolom versi baru.
\item Blok OTORISASI wajib lengkap (Dosen Pengembang RPS, Koordinator RMK, Ka Prodi) sebelum RPS dianggap sah dan dapat diunggah ke SIAKAD.
\item Hal yang belum diatur dalam petunjuk teknis ini diputuskan oleh Ketua Jurusan dan Ketua Program Studi dengan mengacu pada peraturan Politeknik Negeri Malang yang berlaku.
\end{enumerate}

\vspace{1cm}

\begin{flushright}
Malang, \underline{\hspace{3.5cm}}\\
Ketua Jurusan\\[2.2cm]
\textbf{\KetuaJurusan}\\
NIP. \KetuaJurusanNIP
\end{flushright}

\end{document}
